\begin{abstract}
Execution-based evaluation of neurosymbolic reasoning conflates faithful derivation with structural shortcuts: a formalization that embeds the conclusion among its premises proves just as reliably as a genuine one, and reinforcement learning (RL) optimizers exploit this indifference.
We apply leave-one-out ablation to GRPO-trained formalizations, \emph{Solver Trace Ablation} (STA), testing whether each component is necessary for proof success, and characterize structural degeneration under six reward conditions via the resulting \emph{Intensional Structural Reward} (ISR).
ISR reveals that prove-rate gains conceal a compositional shift toward hacking modes, but ISR occupies a paradoxical position: the same leave-one-out structure that makes it informative also makes it gameable.
Tautological formalizations satisfy the necessity criterion by construction, so when ISR is repurposed as a training reward, the model learns to game the diagnostic rather than the task.
In our Qwen3-4B experiments, the resulting Goodhart effect is seed-dependent (3 of 5 seeds converge to gaming, demonstrating that the Goodhart effect can occur under standard training conditions; 2 of 5 retain genuine reasoning), but the \emph{observe-vs-optimize asymmetry} persists across logic types, confirmed by cross-logic (Prolog) replication, while cross-dataset (ProofFOL) evidence reveals dataset-dependent boundary conditions on the effect: structural proxies for formalization integrity are easier to observe than to optimize.\end{abstract}
